\scanpage{59}

\uline{b) réponse à 19.5.6. (iii)}

L'implication $b)\Rightarrow c)$ de 19.5.5 est vraie sans qu'on doive faire aucune hypothèse de séparation. Pour le voir, il suffit de lire le cas $n=1$ de la démonstration ``à la Schwartz''.

Vu l'urgence que tu m'as signalée, et bien qu'il soit peut-être trop tard, je joins une rédaction de la démonstration (copiée sur la primitive)

--- Je n'ai pas encore compris 19.7.

--- J'ai beaucoup aimé le projet EGA I \S\,0 jusqu'à la pg 10, où la topologie du spectre tombe du ciel. Ne faudrait-il pas parler de topologie avant de parler d'espace topologique ?

--- Je ne suis pas capable de rédiger une théorie naïve utilisable d\deleted{\unsure{'}}es intersections, et n'ai pas d'autre remarque à faire sur le ``plan des EGA'' pour le moment

Bien à toi\\

P. Deligne\\
% archive sheet number: 58
